Teorie množin, či její nejrozšířenější formulace,
\href{https://en.wikipedia.org/wiki/Zermelo%E2%80%93Fraenkel_set_theory}{Zermelova-Fraenkelova
teorie množin s axiomem výběru}, je zatím snad nejzdařilejší (na každý pád
nejpoužívanější) pokus o formalisaci základů, na nichž stojí matematika. Čili,
jde o \hyperref[def:teorie]{teorii}, jejímž \hyperref[def:model]{modelem} je
matematika.

Základním pojmem teorie množin, je -- logicky -- \emph{množina}. O
\emph{množinách} můžeme přemýšlet jako o souborech věcí nebo též celcích
složených z částí. Těmto věcem či částem uvnitř množin říkáme jejich
\emph{prvky}. Důležité je však zdůraznit, že \textbf{prvky množin jsou opět
množiny}: žádné jiné \uv{věci} kromě množin v teorii množin nefigurují. Takový
přístup má rozumný filosofický podklad -- je vskutku obtížné najít věc, která se
neskládá z částí. Totiž, vezměme například množinu všech lidí, tj. množinu
jejímiž prvky jsou lidé. Každý člověk je sám o sobě množinou, skládá se mimo
jiné ze svých kostí a svalů. Sama kost i sval jsou množinami -- kosti se
skládají z kolagenu a anorganických materiálů, svaly z měkké tkáně. Takto bychom
mohli pokračovat teoreticky do nekonečna.

\hyperref[def:jazyk]{Jazyk} teorie množin rozšiřuje predikátovou logiku primárně
o pojem \emph{množiny}, symbol náležení $ \in $ a symbol rovnosti $=$. Jsou-li
$x, y$ dvě množiny, pak $x \in y$ je \hyperref[def:vyrok]{výrok} v
\hyperref[def:predikatova-logika]{predikátové logice}, který čteme jako: \uv{$x$
náleží $y$,} též: \uv{$x$ leží v $y$,} nebo: \uv{$x$ je prvkem $y$,} a znamená,
jak jistě tušíte, že množina $x$ je prvkem množiny $y$. Symboly $x = y$ jsou
rovněž výrokem, jež čteme přirozeně jako: \uv{$x$ je rovno $y$.}

Teorie množin s sebou rovněž nese několik dalších užitečných symbolů pro
množinové operace, které si však představíme až spolu s \emph{axiomy}
zaručujícími jejich existenci.

\section{Axiomy teorie množin}
\label{sec:axiomy-teorie-mnozin}
Připomínáme, že slovo \emph{axiom} je označením pro gramaticky správnou
posloupnost symbolů v daném \hyperref[def:jazyk]{jazyce}, která tvoří část
\hyperref[def:teorie]{teorie}; je tedy vlastně \uv{označena za
\emph{pravdivou}}. Protože \hyperref[def:jazyk]{jazyk} teorie množin dědí jazyk
\hyperref[def:predikatova-logika]{predikátové logiky}, jsou axiomy teorie množin
\hyperref[def:vyrok]{výroky}, o kterých prohlásíme, že platí.

Zásluhou mnoha matematiků (převážně) 20. století, je teorie množin redukována na
rozumně malé množství (skupin) axiomů. Cílem však často není získat tu nejmenší
možnou teorii, ale též zachovat jistou \uv{intuitivnost} zvolených axiomů.
Teorie množin nejsouc sic tak malá, jak by jistě mohla být, zachovává si jistý
vyšší stupeň srozumitelnosti.

Probereme se nyní jejími axiomy -- dá se říci od nejlehčích po nejtěžší -- a u
každého předneseme vysvětlení, příměry a snahy o \uv{polidštění}. Mnoho axiomů
též dává možnost vzniknout operacím na množinách, jež nám přijde vhodné uvést
společně s nimi.

\subsection{Axiom existence}
\label{ssec:axiom-existence}

Mnozí teologové v průběhu tisíciletí našeho letopočtu vnímali Boha jako prapůvod
existence -- tvůrce, jenž stvořil z ničeho něco. Otázka, proč ve vesmíru vůbec
něco je, zůstává ve fyzice a filosofii víceméně otevřená. Teorie množin, jejímiž
architekty jsou bezesporu lidé, pročež obsahuje \emph{axiom existence}, který
zaručuje, že v matematickém vesmíru vůbec něco je.

Formulovat jej není pravda obtížné, ale vyžaduje chvíli zamyšlení. Totiž
\hyperref[def:existencni-kvantifikator]{existenční kvantifikátor} umožňuje
vyjádřit, že \emph{existuje $x$}. Avšak, jaké $x$? Podle pravidel
\hyperref[def:predikatova-logika]{predikátové logiky} musí kvantifikátoru
následovat výrok. Jelikož zatím žádné množiny k disposici nemáme, nabízí se
výrok $x = x$. Tedy, axiom existence vlastně říká, že \textbf{existuje množina
rovna sobě samé}. Jeho přesná formulace jest
\begin{equation}
  \label{eq:axiom-existence}
  ( \exists x)(x = x).
\end{equation}
Axiom existence si můžeme pamatovat pod triviálním heslem:
\begin{highlightBox}
  \alert{Existuje něco.}
\end{highlightBox}

\subsection{Axiom extensionality}
\label{ssec:axiom-extensionality}

Název tohoto axiomu je záhadný, ale vyjadřuje přirozenou intuici, že dvě
množiny, které mají přesně ty samé prvky, už nutně musejí být stejné. Heslovitě,
v teorii množin platí:
\begin{highlightBox}
  \alert{Celek je přesně (tedy ani víc ani méně) než součet svých částí.}
\end{highlightBox}
Formulovat jej opět není obtížné, stačí přeformulovat větu: \uv{Množiny $x$ a
$y$ mají stejné prvky.} Inu, to se stane přesně tehdy, když pro každou množinu
platí, že je prvkem $x$ právě tehdy, když je prvkem $y$. V jazyce teorie množin,
výrok
\[
  ( \forall z)(z \in x \Leftrightarrow z \in y)
\]
přesně vyjadřuje, že množiny $x$ a $y$ mají totožné prvky. \emph{Axiom
extensionality} tvrdí, že v~tomto případě již je $x$ rovna $y$. Ještě jednou
vyslovíme \emph{axiom extensionality} slovně a pak jej vyjádříme symbolicky.
Tedy, v teorii množin platí výrok: \uv{Když je každá množina prvkem $x$ právě
tehdy, když je prvkem $y$, pak jsou množiny $x$ a $y$ stejné}. Symbolicky:
\begin{equation}
  \label{eq:axiom-extensionality}
  ( \forall z)(z \in x \Leftrightarrow z \in y) \Rightarrow (x = y).
\end{equation}
\begin{notation}
  Fakt, že množina $a$ má za prvky množiny $b, c$ a $d$ a \emph{žádné jiné}
  zapisujeme
  \[
    a = \{b, c, d\}.
  \]
\end{notation}
\begin{warning}
  Pozor! \hyperref[ssec:axiom-extensionality]{Axiom extensionality} v žádném
  smyslu nehovoří o \emph{uspořádání} nebo \emph{pořadí} prvků v množinách.
  Obecné \textbf{množiny nemají uspořádané prvky}. Vskutku, podle
  \hyperref[ssec:axiom-extensionality]{axiomu extensionality} jsou množiny
  \[
    a = \{x, y, z\} \quad \text{a} \quad b = \{z, y, x\}
  \]
  \emph{naprosto stejné}, přestože jsme se rozhodli napsat prvky množiny $b$ v
  jiném pořadí než prvky množiny $a$.
\end{warning}

Když na levé straně výroku~\eqref{eq:axiom-extensionality} platí pouze implikace
zleva doprava, tj. když platí výrok
\[
  ( \forall z)(z \in x \Rightarrow z \in y),
\]
který lze přetlumočit na: \uv{Každá množina, jež je prvkem $x$, je taky prvkem
$y$,} pak říkáme, že $x$ \emph{je podmnožina} $y$; tento vztah jednoduše
znamená, že $x$ se skládá z některých částí $y$.

\begin{definition}[Podmnožina]
  \label{def:podmnozina}
  O množině $x$ řekneme, že je \emph{podmnožinou} množiny $y$ -- a píšeme $x
  \subseteq y$ -- pokud platí
  \begin{equation}
    \label{eq:podmnozina}
    ( \forall z)(z \in x \Rightarrow z \in y).
  \end{equation}
\end{definition}
\begin{example}
  \label{exam:podmnozina}
  Je-li $x$ množina lidí a $y$ množina \emph{dospělých} lidí, pak platí $y
  \subseteq x$, tedy $y$ je \hyperref[def:podmnozina]{podmnožinou} $x$. Je tomu
  tak pro to, že každý dospělý člověk je samozřejmě též člověk, ale ne každý
  člověk je nezbytně dospělý.
\end{example}

